\section*{Bab \ref{ch:b2:asexual-reproduction} --- Perkembangbiakan Aseksual dan Pengklonaan pada Tumbuhan}

\begin{solution}{exo:b2:asexual-reproduction:1}
\omterm{def:b2:asexual-reproduction:clone}{Klon}: keturunan satu tetua yang serupa secara genetik lewat mitosis.
\omterm{def:b2:asexual-reproduction:clone}{Genet}: individu genetiknya, seluruh yang tumbuh dari satu zigot.
\omterm{def:b2:asexual-reproduction:clone}{Ramet}: satu satuan sebuah \omterm{def:b2:asexual-reproduction:clone}{genet} yang merdeka secara fisiologi.
Organnya: sulur (arbei), rimpang (rumput teki), \omterm{prop:b2:asexual-reproduction:organs}{umbi batang} (kentang),
\omterm{prop:b2:asexual-reproduction:organs}{umbi lapis} (bawang), tunas akar (aspen), tumbuhan kecil daun
(\emph{Kalanchoe}).
\end{solution}

\begin{solution}{exo:b2:asexual-reproduction:2}
Kambium pembuluh batang atas dan \omterm{prop:b2:asexual-reproduction:horticulture}{batang bawahnya} harus sejajar dan
bersentuhan; kalusnya menjembatani potongannya lalu
mendiferensiasikan xilem dan floem menyeberanginya. \omterm{prop:b2:asexual-reproduction:horticulture}{Batang bawahnya}
memasok akar (air, mineral, kekuatan, ketahanan tanah, kendali ukuran),
batang atasnya memasok pucuk berbuahnya (varietasnya). Tiap bagiannya
mempertahankan genomnya sendiri dan selnya tidak pernah melebur: dua
\omterm{def:b2:meiosis-heredity:vocabulary}{genotipe} berdampingan, yaitu sebuah kimera; sebuah hibrid akan
mengusung campuran kedua genomnya di dalam tiap selnya.
\end{solution}

\begin{solution}{exo:b2:asexual-reproduction:3}
\omterm{def:b2:asexual-reproduction:clone}{Pelipatgandaan vegetatif}: sebagian tubuhnya berakar menjadi tumbuhan
baru. \omterm{def:b2:asexual-reproduction:apomixis}{Apomiksis}: sebutir \omterm{def:b2:angiosperm-reproduction:seed}{biji} terbentuk tanpa pembuahan, dan lembaganya
adalah \omterm{def:b2:asexual-reproduction:clone}{klon} induknya. \omterm{def:b2:asexual-reproduction:apomixis}{Partenogenesis}: seekor hewan berkembang dari
telur yang tak dibuahi. Hanya \omterm{def:b2:asexual-reproduction:apomixis}{apomiksis} yang mempertahankan \omterm{def:b2:angiosperm-reproduction:seed}{bijinya} ---
beserta kulitnya, \omterm{def:b2:angiosperm-reproduction:seed}{dormansinya}, dan penyebarannya.
\end{solution}

\begin{solution}{exo:b2:asexual-reproduction:4}
Sterilkan eksplannya; letakkan pada agar berisi gula, mineral, vitamin,
serta auksin dan sitokinin yang seimbang (kalus); naikkan rasio
sitokinin terhadap auksin untuk mengimbas pucuknya; subkulturkan
pucuknya tiap beberapa pekan; pindahkan ke medium kaya auksin untuk
mengakarkannya; sapihlah tumbuhan kecilnya di bawah kelembapan tinggi
lalu di dalam tanah.
\end{solution}

\begin{solution}{exo:b2:asexual-reproduction:5}
Tiap tumbuhannya digantikan $1 + 24 = 25$ tumbuhan: sesudah 4 tahun $25^{4} =
\num{390625}$ tumbuhan, yang membutuhkan \qty{39000}{m^2},
hampir empat hektar. Ruang, cahaya, dan haranya, jangkauan sulurnya yang
pendek, serta kematian tumbuhan karena berdesakan menghentikan fase
eksponensialnya dalam beberapa tahun; \omterm{def:b2:asexual-reproduction:clone}{klonnya} lalu maju sebagai sebuah
muka.
\end{solution}

\begin{solution}{exo:b2:asexual-reproduction:6}
$\num{14000}/130 \approx 110$ generasi batang. Jari-jari sebuah cakram
seluas \qty{4.3e5}{m^2}: $\sqrt{4.3\times 10^{5}/\pi} = \qty{370}{m}$;
kemajuannya $370/\num{14000} = \qty{0.026}{m}$ tiap tahun, yaitu kurang
dari tiga sentimeter.
\end{solution}

\begin{solution}{exo:b2:asexual-reproduction:7}
$p_t = 2^{t}p_0/(1 + (2^{t} - 1)p_0)$: $t = 5$: $0.0032$; $t = 10$:
$0.1024/1.1023 = 0.093$; $t = 15$: $3.28/4.28 = 0.77$. Ia tidak pernah
menyebar kalau kebugarannya di bawah separuh kebugaran garis keturunan
seksualnya: kebugaran nisbi $w$ memberi $q' = 2wq$, yang meluruh bagi $w <
\tfrac12$.
\end{solution}

\begin{solution}{exo:b2:asexual-reproduction:8}
$10 : 1$: akar; $1 : 1$: kalus yang tak terdiferensiasi; $1 : 10$:
pucuk. Sebuah setek sudah memiliki pucuk dan membutuhkan akar, yang
diimbas auksin pada pangkal potongannya.
\end{solution}

\begin{solution}{exo:b2:asexual-reproduction:9}
1, 10, 100, \num{1000}, \num{10000} kg: sesudah musim ketiga baru
\qty{1}{t}, sesudah musim keempat \qty{10}{t}: empat musim. \omterm{def:b2:angiosperm-reproduction:seed}{Biji} sejati
memberi ribuan \omterm{def:b2:angiosperm-reproduction:seed}{biji} tiap tumbuhan tetapi kentang adalah tetraploid yang
\omterm{def:b2:meiosis-heredity:vocabulary}{heterozigot}, sehingga kecambahnya memilah dan tak satu pun merupakan
varietasnya.
\end{solution}

\begin{solution}{exo:b2:asexual-reproduction:10}
$e^{-5} = 0.0067$: 670 individu bebas \omterm{def:b2:genome-diversification:mutation}{mutasi} pada populasi yang besar,
kira-kira 7 pada yang kecil. Tujuh individu semuanya dapat gagal
meninggalkan keturunan secara kebetulan dalam beberapa generasi
(peluang bahwa tak satu pun berbiak seorde $e^{-7}$ tiap generasi), dan
sesudah itu golongan terbaiknya lenyap selamanya dan golongan
terbaik berikutnya menjadi yang terbaik: satu ketukan ratchetnya, yang
berulang sampai bebannya tak tertanggungkan. Dengan 670, kehilangannya
praktis mustahil --- sampai sebuah leher botol memerkecil $N$ atau \omterm{thm:b2:genome-diversification:rate}{laju
mutasinya} naik; ratchetnya hanya pernah berputar satu arah.
\end{solution}

\begin{solution}{exo:b2:asexual-reproduction:11}
Di dalam \omterm{def:b2:asexual-reproduction:clone}{klonnya} tiap inangnya peka: galurnya menyebar tanpa terkendali
pada laju sporanya mengembara. Di dalam populasi seksualnya, \omterm{def:b2:meiosis-heredity:vocabulary}{genotipe}
yang dikalahkannya adalah satu dari $2^{10} = 1024$ gabungan, yang
hadir pada kira-kira seperseribu tumbuhannya, terserak di antara
tetangga yang kebal, sehingga wabahnya kelaparan. Lumper adalah satu
\omterm{def:b2:asexual-reproduction:clone}{klon} yang ditanam di sebagian besar sebuah negeri; hawarnya mendapati
tiap tumbuhannya terbuka, dan tidak ada keragaman untuk dimuliakan
sampai varietas lain diimpor.
\end{solution}

\begin{solution}{exo:b2:asexual-reproduction:12}
\omterm{prop:b2:asexual-reproduction:whysex}{Ratchet Muller}, penggabungan \omterm{def:b2:genome-diversification:mutation}{mutasi} yang menguntungkan, dan Ratu Merah
masing-masing memberi seks sebuah manfaat yang dapat melampaui faktor
dua di dunia yang besar, terparasit, dan berubah. Rotifera bdeloid
adalah kasus yang canggung: empat puluh juta tahun tanpa jantan dan
tanpa keruntuhan ratchet. Mereka tampaknya lolos lewat ketahanan
pengeringan yang ekstrem (yang membunuh parasitnya), lewat penyerapan
DNA asing, dan lewat \omterm{prop:b2:genome-diversification:repair}{perbaikan DNA} yang luar biasa berdaya guna ---
kekecualian yang menunjukkan apa yang biasanya dituntut kaidahnya.
\end{solution}

\begin{solution}{pb:b2:asexual-reproduction:1}
\textbf{1.} $4\times 10^{4}$ tumbuhan.
\textbf{2.} $1 + 6\times 2 = 13$.
\textbf{3.} \num{1300}; \num{16900}; \num{219700}.
\textbf{4.} $13^{t} = 400$: $t = \ln 400/\ln 13 = 5.99/2.56 = 2.3$
tahun --- penuh selama tahun ketiga.
\textbf{5.} Jari-jari cakram seluas satu hektar
$\sqrt{10^{4}/\pi} = \qty{56}{m}$; pada \qty{0.5}{m} tiap tahun,
\num{113} tahun.
\textbf{6.} Seratus pendiri yang terserak memenuhi ladangnya secara
eksponensial dari seratus muka dalam dua tahun; satu pendiri akan
memenuhi lingkungannya dalam dua tahun lalu merangkak selama seabad.
\textbf{7.} $5^{n} \ge 10^{6}$: $n \ge 8.6$, yaitu sembilan subkultur
($1.95\times 10^{6}$ pucuk), 36~pekan.
\textbf{8.} $0.8\times 1.95\times 10^{6} = 1.6\times 10^{6}$ tumbuhan.
\textbf{9.} $0.98$; $0.98^{10} = 0.82$.
\textbf{10.} $10\times 0.98 = 9.8$ galur bersih yang diharapkan; dari
sepuluh setek, satu.
\textbf{11.} Virusnya mengembara lewat plasmodesmata dan floemnya, dan
kubah pucuknya tidak memiliki hubungan pembuluh serta membelah lebih
cepat daripada penyebaran virusnya, sehingga sepersepuluh milimeter
terluarnya biasanya tak terjangkiti.
\textbf{12.} Laboratoriumnya: tumbuhan yang seragam dan bebas virus
seketika, tetapi mahal, dan tiap \omterm{def:b2:genome-diversification:mutation}{mutasi} atau pencemaran di dalam
kulturnya dilipatgandakan sejuta kali. Sulurnya: cuma-cuma dan di
tempat, tetapi tiga musim hilang dan tiap virus di dalam tumbuhan
induknya terbawa serta.
\textbf{13.} Kaum apomiknya berjumlah $pN$ dan memberi $F$ \omterm{def:b2:angiosperm-reproduction:seed}{biji}
masing-masing; kaum seksualnya $(1 - p)N$ memberi $F/2$ \omterm{def:b2:angiosperm-reproduction:seed}{biji}
masing-masing (separuh upayanya adalah \omterm{def:b2:plant-life-cycles:heterospory}{serbuk sari}, yang tidak membuat
\omterm{def:b2:angiosperm-reproduction:seed}{biji}); bagian apomik atas \omterm{def:b2:angiosperm-reproduction:seed}{bijinya} adalah $Fp/(Fp + F(1 - p)/2)
= 2p/(1 + p)$.
\textbf{14.} $p_5 = 0.032/1.031 = 0.031$; $p_{10} = 1.024/2.023 =
0.51$: ia melewati separuh pada generasi ke-10.
\textbf{15.} $p' = 2p(1 - cp)/\bigl(2p(1 - cp) + (1 - p)\bigr)$.
\textbf{16.} $p' = p$ ketika $2(1 - cp^*) = 1$: $p^* = 1/2c = 0.625$.
Di bawah $p^*$, $2(1 - cp) > 1$ dan $p$ naik; di atasnya, ia turun:
mantap.
\textbf{17.} $p^* \ge 1$ ketika $c \le \tfrac12$: apomiknya mengambil
alih kecuali kalau parasitnya memakan lebih dari separuh keluarannya
ketika ia lazim.
\textbf{18.} $c = 0.6$: $p^* = 0.83$. Untuk mempertahankan seks di
dalam populasinya, parasitnya harus, pada frekuensi apomik yang tinggi,
memangkas kebugaran \omterm{def:b2:asexual-reproduction:clone}{klonnya} lebih besar daripada keunggulan dua kali
lipatnya --- Ratu Merah harus berlari cepat.
\textbf{19.} Tiga perangkat kromosom tidak dapat berpasangan pada
\omterm{def:b2:meiosis-heredity:meiosis}{meiosis}, sehingga gametnya tak berimbang; ia terkunci di luar
rekombinasi, dan \omterm{prop:b2:asexual-reproduction:whysex}{ratchet Muller} serta parasitnya bekerja padanya tanpa
perlawanan --- itulah sebabnya garis keturunan rambutan kuning apomik
masih muda dan terus dibuat ulang dari moyang seksualnya.
\textbf{20.} $10^{-9}\times 4.8\times 10^{8} = 0.48$ tiap pembelahan.
\textbf{21.} $2\times \num{14000}\times 20 = 5.6\times 10^{5}$
pembelahan; $0.48\times 5.6\times 10^{5} = 2.7\times 10^{5}$
\omterm{def:b2:genome-diversification:mutation}{mutasi}.
\textbf{22.} $2.7\times 10^{5}/4.8\times 10^{8} = 5.6\times 10^{-4}$
genomnya; kira-kira \num{5400} perbedaan penyandi.
\textbf{23.} Persaingan antargaris keturunan sel di dalam meristemnya
(sebuah galur mutan yang membelah lebih lambat tersingkir) dan
antarramet (sebuah batang yang mengusung \omterm{def:b2:genome-diversification:mutation}{mutasi} yang buruk bertunas
lebih sedikit). Lebih lemah karena sebagian besar \omterm{def:b2:genome-diversification:mutation}{mutasi} somanya
\omterm{def:b2:meiosis-heredity:vocabulary}{heterozigot} dan tertutupi, dan tidak ada yang menyingkapkannya: tidak
ada \omterm{def:b2:meiosis-heredity:meiosis}{meiosis} yang menjadikannya \omterm{def:b2:meiosis-heredity:vocabulary}{homozigot} dan tidak ada tahap zigot
\omterm{def:b2:unicellular-diversity:unicellular}{bersel tunggal} yang menguji seluruh genomnya di dalam satu sel.
\textbf{24.} Tidak seragam: batangnya berbeda pada ratusan ribu tapak.
Tetapi perbedaannya sebagian besar netral dan tertutupi, sehingga
\omterm{def:b2:asexual-reproduction:clone}{klonnya} secara efektif satu \omterm{def:b2:meiosis-heredity:vocabulary}{genotipe} --- sebuah mosaik varian soma satu
\omterm{def:b2:asexual-reproduction:clone}{genet}.
\textbf{25.} Ladangnya penuh sesudah 2,3 tahun; satu meristem memberi
$1.6\times 10^{6}$ tumbuhan dalam sembilan bulan; dengan $c = 0.8$
apomiknya menetap pada $p^* = 0.625$, dan kaum seksualnya
mempertahankan \qty{37.5}{\%}.
\end{solution}
